\documentclass[11pt,a4paper]{article}

% =========================================================
% PACOTES
% =========================================================
\usepackage[utf8]{inputenc}
\usepackage[T1]{fontenc}
\usepackage[portuguese]{babel}
\usepackage{lmodern}

\usepackage[a4paper,
top=1.8cm,
bottom=1.8cm,
left=1.8cm,
right=1.8cm]{geometry}

\usepackage{amsmath,amssymb}
\usepackage{array}
\usepackage{booktabs}
\usepackage{tabularx}
\usepackage{multirow}
\usepackage{multicol}
\usepackage{enumitem}
\usepackage{siunitx}
\usepackage{xcolor}
\usepackage{graphicx}
\usepackage{float}

\usepackage{tikz}
\usetikzlibrary{
	arrows.meta,
	positioning,
	calc,
	decorations.pathmorphing,
	shapes.geometric,
	backgrounds
}

\usepackage{pgfplots}
\pgfplotsset{compat=1.18}

\usepackage[most]{tcolorbox}
\usepackage{fancyhdr}
\usepackage{lastpage}
\usepackage{hyperref}

% =========================================================
% CORES
% =========================================================
\definecolor{eublue}{HTML}{124E78}
\definecolor{eulightblue}{HTML}{D9EEF7}
\definecolor{eudark}{HTML}{16324F}
\definecolor{eugreen}{HTML}{3A7D44}
\definecolor{eulightgreen}{HTML}{E4F2E7}
\definecolor{euorange}{HTML}{E58E26}
\definecolor{eulightorange}{HTML}{FCEBD8}
\definecolor{eugray}{HTML}{F2F4F5}
\definecolor{eured}{HTML}{B03A2E}
\definecolor{eupurple}{HTML}{6C5B7B}

% =========================================================
% CONFIGURAÇÕES
% =========================================================
\sisetup{
	output-decimal-marker={,},
	per-mode=symbol
}

\hypersetup{
	colorlinks=true,
	linkcolor=eublue,
	urlcolor=eublue
}

\setlength{\parindent}{0pt}
\setlength{\parskip}{4pt}

\renewcommand{\arraystretch}{1.35}

% =========================================================
% CABEÇALHO E RODAPÉ
% =========================================================
\pagestyle{fancy}
\fancyhf{}

\fancyhead[L]{\textbf{\textcolor{eublue}{EUFISICA exercícios}}}
\fancyhead[R]{\textcolor{eudark}{Estrutura Atómica e Nanotecnologia}}

\fancyfoot[L]{Prof. José Gonçalves}
\fancyfoot[C]{\texttt{EUFISICA}}
\fancyfoot[R]{Página \thepage\ de \pageref{LastPage}}

\renewcommand{\headrulewidth}{0.6pt}
\renewcommand{\footrulewidth}{0.4pt}

% =========================================================
% CAIXAS
% =========================================================
\newcounter{exercicio}

\newenvironment{exercicio}[1]
{
	\refstepcounter{exercicio}
	\begin{tcolorbox}[
		enhanced,
		breakable,
		colback=white,
		colframe=eublue,
		boxrule=0.8pt,
		arc=2mm,
		left=3mm,
		right=3mm,
		top=2mm,
		bottom=2mm,
		title={
			\textbf{Exercício \theexercicio}
			\hspace{1mm}--\hspace{1mm}
			#1
		},
		coltitle=white,
		colbacktitle=eublue,
		fonttitle=\bfseries
		]
	}
	{
	\end{tcolorbox}
}

\newtcolorbox{recorda}{
	enhanced,
	colback=eulightblue,
	colframe=eublue,
	title=\textbf{Recorda},
	coltitle=eudark,
	colbacktitle=eulightblue,
	fonttitle=\bfseries,
	arc=2mm
}

\newtcolorbox{stem}{
	enhanced,
	colback=eulightgreen,
	colframe=eugreen,
	title=\textbf{STEM | Ciência, Tecnologia e Mundo Real},
	coltitle=eugreen,
	colbacktitle=eulightgreen,
	fonttitle=\bfseries,
	arc=2mm
}

\newtcolorbox{desafio}{
	enhanced,
	colback=eulightorange,
	colframe=euorange,
	title=\textbf{Desafio},
	coltitle=euorange,
	colbacktitle=eulightorange,
	fonttitle=\bfseries,
	arc=2mm
}

% =========================================================
% LINHA DE RESPOSTA
% =========================================================
\newcommand{\linha}{
	\vspace{2mm}
	\noindent\makebox[\linewidth]{\dotfill}
}

\newcommand{\linhas}[1]{
	\par
	\foreach \x in {1,...,#1}{
		\vspace{3mm}
		\noindent\makebox[\linewidth]{\dotfill}\par
	}
}

% =========================================================
% DOCUMENTO
% =========================================================
\begin{document}
	
	% =========================================================
	% TÍTULO
	% =========================================================
	
	\begin{tcolorbox}[
		enhanced,
		colback=eudark,
		colframe=eudark,
		arc=3mm,
		boxrule=0pt,
		left=5mm,
		right=5mm,
		top=5mm,
		bottom=5mm
		]
		\begin{center}
			
			{\Huge\bfseries\color{white} EUFISICA exercícios}
			
			\vspace{2mm}
			
			{\Large\bfseries\color{white}
				Estrutura Atómica e Nanotecnologia}
			
			\vspace{3mm}
			
			{\large\color{eulightblue}
				Número atómico • Número de massa • Massa atómica relativa}
			
			{\large\color{eulightblue}
				Protões • Neutrões • Eletrões • Nanotecnologia}
			
			\vspace{3mm}
			
			{\color{white}\textbf{Prof. José Gonçalves}}
			
		\end{center}
	\end{tcolorbox}
	
	\vspace{3mm}
	
	\begin{tabularx}{\textwidth}{|X|X|X|}
		\hline
		\textbf{Nome:} & \textbf{Turma:} & \textbf{Data:} \\[6mm]
		\hline
	\end{tabularx}
	
	\vspace{5mm}
	
	% =========================================================
	% OBJETIVOS
	% =========================================================
	
	\begin{tcolorbox}[
		colback=eugray,
		colframe=eudark,
		title=\textbf{Objetivos de aprendizagem}
		]
		No final desta ficha deverás ser capaz de:
		
		\begin{itemize}[leftmargin=6mm,itemsep=1mm]
			\item distinguir \textbf{número atómico} de \textbf{número de massa};
			\item determinar o número de \textbf{protões, neutrões e eletrões};
			\item interpretar corretamente a representação simbólica de um átomo ou ião;
			\item calcular a \textbf{massa atómica relativa} a partir das abundâncias isotópicas;
			\item relacionar a escala nanométrica com aplicações atuais da \textbf{nanotecnologia};
			\item compreender a importância da relação entre \textbf{área superficial e volume} à escala nano.
		\end{itemize}
	\end{tcolorbox}
	
	% =========================================================
	% SECÇÃO 1
	% =========================================================
	
	\section*{\textcolor{eublue}{1. Número atómico e número de massa}}
	
	\begin{recorda}
		
		Um átomo pode representar-se simbolicamente por:
		
		\[
		{}^{A}_{Z}\mathrm{X}
		\]
		
		onde:
		
		\[
		\boxed{Z = \text{número atómico}}
		\]
		
		\[
		\boxed{A = \text{número de massa}}
		\]
		
		O número atómico, \(Z\), corresponde ao número de \textbf{protões} do núcleo.
		
		O número de massa, \(A\), corresponde à soma do número de \textbf{protões} com o número de \textbf{neutrões}:
		
		\[
		\boxed{A=Z+N}
		\]
		
		sendo \(N\) o número de neutrões.
		
		Consequentemente,
		
		\[
		\boxed{N=A-Z}
		\]
		
	\end{recorda}
	
	% ---------------------------------------------------------
	% IMAGEM REPRESENTAÇÃO DO NÚCLIDO
	% ---------------------------------------------------------
	
	\begin{figure}[H]
		\centering
		
		\begin{tikzpicture}[scale=1.15]
			
			\node[font=\fontsize{42}{42}\selectfont\bfseries,text=eudark]
			(X) at (0,0) {X};
			
			\node[font=\Large\bfseries,text=eublue]
			at (-0.85,0.65) {\(A\)};
			
			\node[font=\Large\bfseries,text=eugreen]
			at (-0.85,-0.55) {\(Z\)};
			
			\draw[-{Latex[length=3mm]},thick,eublue]
			(-1.25,1.4) -- (-0.9,0.85);
			
			\node[anchor=east,text=eublue,font=\bfseries]
			at (-1.3,1.4) {Número de massa};
			
			\draw[-{Latex[length=3mm]},thick,eugreen]
			(-1.25,-1.25) -- (-0.9,-0.7);
			
			\node[anchor=east,text=eugreen,font=\bfseries]
			at (-1.3,-1.25) {Número atómico};
			
			\draw[-{Latex[length=3mm]},thick,eudark]
			(2.6,0) -- (0.65,0);
			
			\node[anchor=west,text=eudark,font=\bfseries]
			at (2.7,0) {Símbolo químico};
			
		\end{tikzpicture}
		
		\caption{Representação simbólica geral de um nuclídeo.}
	\end{figure}
	
	% =========================================================
	
	\begin{exercicio}{Identificar \(Z\), \(A\), protões e neutrões}
		
		Considera o átomo representado por:
		
		\[
		{}^{23}_{11}\mathrm{Na}
		\]
		
		Completa:
		
		\begin{center}
			\begin{tabular}{|l|c|}
				\hline
				\textbf{Grandeza} & \textbf{Valor} \\ \hline
				Número atómico, \(Z\) & \\[4mm] \hline
				Número de massa, \(A\) & \\[4mm] \hline
				Número de protões & \\[4mm] \hline
				Número de neutrões & \\[4mm] \hline
				Número de eletrões, se o átomo for neutro & \\[4mm] \hline
			\end{tabular}
		\end{center}
		
	\end{exercicio}
	
	% =========================================================
	
	\begin{exercicio}{Completar representações simbólicas}
		
		Completa a tabela.
		
		\begin{center}
			\small
			
			\begin{tabular}{|c|c|c|c|c|c|}
				\hline
				\textbf{Espécie}
				&
				\(\mathbf{Z}\)
				&
				\(\mathbf{A}\)
				&
				\textbf{Protões}
				&
				\textbf{Neutrões}
				&
				\textbf{Eletrões}
				\\ \hline
				
				\({}^{16}_{8}\mathrm{O}\)
				& & & & & \\[4mm] \hline
				
				\({}^{27}_{13}\mathrm{Al}\)
				& & & & & \\[4mm] \hline
				
				\({}^{40}_{20}\mathrm{Ca}\)
				& & & & & \\[4mm] \hline
				
				\({}^{35}_{17}\mathrm{Cl}\)
				& & & & & \\[4mm] \hline
				
			\end{tabular}
		\end{center}
		
	\end{exercicio}
	
	% =========================================================
	
	\begin{exercicio}{Descobrir o elemento}
		
		Um átomo possui:
		
		\[
		17\text{ protões}
		\qquad\text{e}\qquad
		18\text{ neutrões}.
		\]
		
		\begin{enumerate}[label=\alph*)]
			\item Determina o número atómico.
			\linhas{1}
			
			\item Determina o número de massa.
			\linhas{1}
			
			\item Quantos eletrões possui, se for eletricamente neutro?
			\linhas{1}
			
			\item Sabendo que o elemento de número atómico 17 é o cloro, representa simbolicamente este átomo.
			\linhas{1}
		\end{enumerate}
		
	\end{exercicio}
	
	% =========================================================
	% SECÇÃO 2
	% =========================================================
	
	\section*{\textcolor{eublue}{2. Protões, neutrões e eletrões}}
	
	\begin{recorda}
		
		\begin{center}
			\begin{tabular}{cccc}
				\toprule
				\textbf{Partícula}
				&
				\textbf{Símbolo}
				&
				\textbf{Carga relativa}
				&
				\textbf{Localização}
				\\
				\midrule
				Protão & \(p^+\) & \(+1\) & Núcleo \\
				
				Neutrão & \(n^0\) & \(0\) & Núcleo \\
				
				Eletrão & \(e^-\) & \(-1\) & Nuvem eletrónica \\
				\bottomrule
			\end{tabular}
		\end{center}
		
		Num átomo eletricamente neutro:
		
		\[
		\boxed{n(e^-)=n(p^+)=Z}
		\]
		
	\end{recorda}
	
	% =========================================================
	% MODELO DE ÁTOMO
	% =========================================================
	
	\begin{figure}[H]
		\centering
		
		\begin{tikzpicture}[scale=1.05]
			
			% órbitas
			\draw[eublue!55,thick] (0,0) ellipse (3.3cm and 1.3cm);
			\draw[eublue!55,thick,rotate=60] (0,0) ellipse (3.3cm and 1.3cm);
			\draw[eublue!55,thick,rotate=-60] (0,0) ellipse (3.3cm and 1.3cm);
			
			% núcleo
			\fill[eulightorange] (0,0) circle (0.9cm);
			\draw[euorange,very thick] (0,0) circle (0.9cm);
			
			% nucleões
			\fill[eured!75] (-0.28,0.25) circle (0.27);
			\node[white,font=\scriptsize\bfseries] at (-0.28,0.25) {\(p^+\)};
			
			\fill[eugray] (0.3,0.28) circle (0.27);
			\node[eudark,font=\scriptsize\bfseries] at (0.3,0.28) {\(n^0\)};
			
			\fill[eugray] (-0.25,-0.3) circle (0.27);
			\node[eudark,font=\scriptsize\bfseries] at (-0.25,-0.3) {\(n^0\)};
			
			\fill[eured!75] (0.32,-0.28) circle (0.27);
			\node[white,font=\scriptsize\bfseries] at (0.32,-0.28) {\(p^+\)};
			
			% eletrões
			\fill[eublue] (3.3,0) circle (0.16);
			\node[eublue,font=\bfseries] at (3.65,0.25) {\(e^-\)};
			
			\fill[eublue] (-1.65,2.48) circle (0.16);
			\node[eublue,font=\bfseries] at (-2,2.75) {\(e^-\)};
			
			\fill[eublue] (-1.65,-2.48) circle (0.16);
			\node[eublue,font=\bfseries] at (-2,-2.75) {\(e^-\)};
			
			% legendas
			\draw[-{Latex},thick,eudark] (0.7,0.65) -- (2.0,1.8);
			\node[anchor=west,font=\bfseries] at (2.05,1.85) {Núcleo};
			
			\draw[-{Latex},thick,eublue] (-3.8,1.2) -- (-2.7,0.85);
			\node[anchor=east,font=\bfseries,text=eublue]
			at (-3.85,1.2) {Nuvem eletrónica};
			
		\end{tikzpicture}
		
		\caption{Representação esquemática de um átomo. Não está à escala.}
	\end{figure}
	
	% =========================================================
	
	\begin{exercicio}{Interpretar um modelo atómico}
		
		Observa a figura anterior.
		
		\begin{enumerate}[label=\alph*)]
			\item Que partículas se encontram no núcleo?
			\linhas{1}
			
			\item Qual das partículas possui carga elétrica negativa?
			\linhas{1}
			
			\item Qual das partículas não possui carga elétrica?
			\linhas{1}
			
			\item A representação encontra-se à escala? Justifica.
			\linhas{2}
		\end{enumerate}
		
	\end{exercicio}
	
	% =========================================================
	
	\begin{exercicio}{Átomos e iões}
		
		Completa a tabela seguinte.
		
		\begin{center}
			\small
			
			\begin{tabular}{|c|c|c|c|c|c|c|}
				\hline
				\textbf{Espécie}
				&
				\(\mathbf{Z}\)
				&
				\(\mathbf{A}\)
				&
				\(\mathbf{p^+}\)
				&
				\(\mathbf{n^0}\)
				&
				\(\mathbf{e^-}\)
				&
				\textbf{Carga}
				\\ \hline
				
				\({}^{23}_{11}\mathrm{Na}\)
				& & & & & & \(0\)
				\\[4mm]\hline
				
				\({}^{23}_{11}\mathrm{Na}^{+}\)
				& & & & & & 
				\\[4mm]\hline
				
				\({}^{16}_{8}\mathrm{O}^{2-}\)
				& & & & & &
				\\[4mm]\hline
				
				\({}^{27}_{13}\mathrm{Al}^{3+}\)
				& & & & & &
				\\[4mm]\hline
				
				\({}^{35}_{17}\mathrm{Cl}^{-}\)
				& & & & & &
				\\[4mm]\hline
				
			\end{tabular}
		\end{center}
		
	\end{exercicio}
	
	% =========================================================
	
	\begin{exercicio}{Pensar antes de calcular}
		
		Classifica cada afirmação como \textbf{Verdadeira (V)} ou \textbf{Falsa (F)}.
		Corrige as afirmações falsas.
		
		\begin{enumerate}[label=\alph*)]
			
			\item[\(\square\)] O número atómico corresponde ao número de neutrões.
			
			\item[\(\square\)] O número de massa corresponde ao número total de protões e neutrões.
			
			\item[\(\square\)] Todos os átomos de um mesmo elemento apresentam o mesmo número de protões.
			
			\item[\(\square\)] Um átomo neutro apresenta igual número de protões e eletrões.
			
			\item[\(\square\)] Dois isótopos de um mesmo elemento possuem diferente número atómico.
			
			\item[\(\square\)] Quando um átomo perde um eletrão origina um ião positivo.
			
		\end{enumerate}
		
		\linhas{3}
		
	\end{exercicio}
	
	% =========================================================
	% SECÇÃO 3
	% =========================================================
	
	\section*{\textcolor{eublue}{3. Massa atómica relativa}}
	
	\begin{recorda}
		
		A maioria dos elementos químicos existe na Natureza sob a forma de uma
		\textbf{mistura de isótopos}.
		
		A massa atómica relativa de um elemento, \(A_r\), depende:
		
		\begin{itemize}
			\item das massas dos seus isótopos;
			\item das respetivas abundâncias relativas.
		\end{itemize}
		
		Para dois isótopos:
		
		\[
		\boxed{
			A_r=
			\frac{
				m_1\times \%_1+
				m_2\times \%_2
			}{100}
		}
		\]
		
		A massa atómica relativa é uma \textbf{grandeza adimensional}, isto é,
		não possui unidade.
		
	\end{recorda}
	
	% =========================================================
	
	\begin{exercicio}{Os isótopos do cloro}
		
		Na Natureza, o cloro é constituído principalmente por dois isótopos:
		
		\begin{center}
			\begin{tabular}{ccc}
				\toprule
				\textbf{Isótopo}
				&
				\textbf{Massa isotópica}
				&
				\textbf{Abundância}
				\\
				\midrule
				\({}^{35}\mathrm{Cl}\) & \(34,969\) & \(75,77\,\%\) \\
				
				\({}^{37}\mathrm{Cl}\) & \(36,966\) & \(24,23\,\%\) \\
				\bottomrule
			\end{tabular}
		\end{center}
		
		O gráfico apresenta as abundâncias relativas.
		
		\begin{figure}[H]
			\centering
			
			\begin{tikzpicture}
				
				\begin{axis}[
					width=10.5cm,
					height=6cm,
					ybar,
					ymin=0,
					ymax=100,
					ylabel={Abundância relativa (\%)},
					xlabel={Isótopo},
					symbolic x coords={Cl-35,Cl-37},
					xtick=data,
					bar width=25pt,
					nodes near coords,
					nodes near coords align={vertical},
					axis line style={eudark},
					tick label style={font=\small},
					label style={font=\small\bfseries},
					grid=major,
					grid style={dashed,eugray}
					]
					
					\addplot[
					fill=eublue!70,
					draw=eublue
					]
					coordinates {
						(Cl-35,75.77)
						(Cl-37,24.23)
					};
					
				\end{axis}
			\end{tikzpicture}
			
			\caption{Abundância relativa dos principais isótopos naturais do cloro.}
		\end{figure}
		
		\begin{enumerate}[label=\alph*)]
			
			\item Qual é o isótopo mais abundante?
			\linhas{1}
			
			\item Calcula a massa atómica relativa do cloro.
			
			\[
			A_r(\mathrm{Cl})=
			\]
			
			\linhas{2}
			
			\item O resultado encontra-se mais próximo de 35 ou de 37?
			Explica porquê.
			\linhas{2}
			
		\end{enumerate}
		
	\end{exercicio}
	
	% =========================================================
	
	\begin{exercicio}{Massa atómica relativa do magnésio}
		
		Uma amostra natural de magnésio apresenta aproximadamente a seguinte
		composição:
		
		\begin{center}
			\begin{tabular}{ccc}
				\toprule
				\textbf{Isótopo} & \textbf{Massa aproximada} & \textbf{Abundância} \\
				\midrule
				\({}^{24}\mathrm{Mg}\) & 24 & \(78,99\,\%\)\\
				\({}^{25}\mathrm{Mg}\) & 25 & \(10,00\,\%\)\\
				\({}^{26}\mathrm{Mg}\) & 26 & \(11,01\,\%\)\\
				\bottomrule
			\end{tabular}
		\end{center}
		
		\begin{enumerate}[label=\alph*)]
			\item Calcula uma estimativa para \(A_r(\mathrm{Mg})\).
			\linhas{3}
			
			\item Explica por que razão a massa atómica relativa não é, em geral,
			um número inteiro.
			\linhas{2}
		\end{enumerate}
		
	\end{exercicio}
	
	% =========================================================
	
	\begin{desafio}
		
		Um elemento fictício \(\mathrm{X}\) possui apenas dois isótopos:
		
		\[
		{}^{10}\mathrm{X}
		\qquad\text{e}\qquad
		{}^{11}\mathrm{X}
		\]
		
		Sabendo que:
		
		\[
		A_r(\mathrm{X})=10,8
		\]
		
		determina a percentagem aproximada de cada isótopo na Natureza.
		
		\textbf{Sugestão:} se \(x\) for a percentagem de
		\({}^{10}\mathrm{X}\), então a percentagem de
		\({}^{11}\mathrm{X}\) será \(100-x\).
		
		\linhas{4}
		
	\end{desafio}
	
	% =========================================================
	% NOVA PÁGINA
	% =========================================================
	
	\newpage
	
	\section*{\textcolor{eugreen}{4. Nanotecnologia}}
	
	\begin{stem}
		
		O prefixo \textbf{nano} significa \(10^{-9}\).
		
		Assim,
		
		\[
		\boxed{
			1\,\mathrm{nm}=10^{-9}\,\mathrm{m}
		}
		\]
		
		Um nanómetro corresponde à milionésima parte de um milímetro.
		
		A \textbf{nanotecnologia} envolve a investigação e manipulação da matéria
		a dimensões extremamente pequenas, geralmente da ordem de alguns
		nanómetros até cerca de uma centena de nanómetros.
		
		Nesta escala, os materiais podem apresentar propriedades diferentes das
		que possuem à escala macroscópica.
		
	\end{stem}
	
	% =========================================================
	% ESCALA NANO
	% =========================================================
	
	\begin{figure}[H]
		\centering
		
		\begin{tikzpicture}[x=1.45cm,y=1cm]
			
			% linha
			\draw[-{Latex[length=3mm]},very thick,eudark]
			(0,0)--(8,0);
			
			% marcas
			\foreach \x/\lab in {
				0/{\(10^{-10}\)},
				1.5/{\(10^{-9}\)},
				3/{\(10^{-8}\)},
				4.5/{\(10^{-7}\)},
				6/{\(10^{-6}\)},
				7.5/{\(10^{-5}\)}
			}
			{
				\draw[thick] (\x,0.15)--(\x,-0.15);
				\node[below] at (\x,-0.15) {\lab};
			}
			
			\node[below=8mm,font=\small\bfseries]
			at (3.75,0) {Dimensão aproximada / m};
			
			% objetos
			\node[
			draw=eupurple,
			fill=eupurple!10,
			rounded corners,
			align=center
			] at (0.2,1.0)
			{Átomo\\\(\sim0,1\,\mathrm{nm}\)};
			
			\node[
			draw=eugreen,
			fill=eugreen!10,
			rounded corners,
			align=center
			] at (2.3,1.5)
			{Nanopartícula\\\(\sim10\,\mathrm{nm}\)};
			
			\node[
			draw=eured,
			fill=eured!10,
			rounded corners,
			align=center
			] at (4.5,1.0)
			{Vírus\\\(\sim100\,\mathrm{nm}\)};
			
			\node[
			draw=euorange,
			fill=euorange!10,
			rounded corners,
			align=center
			] at (6.3,1.5)
			{Bactéria\\\(\sim1\,\mathrm{\mu m}\)};
			
			\node[
			draw=eublue,
			fill=eublue!10,
			rounded corners,
			align=center
			] at (7.6,1.0)
			{Cabelo\\\(\sim50{-}100\,\mathrm{\mu m}\)};
			
		\end{tikzpicture}
		
		\caption{Comparação aproximada de diferentes escalas de dimensão.}
	\end{figure}
	
	% =========================================================
	
	\begin{exercicio}{Da escala macroscópica à nanoescala}
		
		Completa as conversões.
		
		\begin{enumerate}[label=\alph*)]
			
			\item
			\[
			1\,\mathrm{nm}
			=
			\underline{\hspace{3cm}}\ \mathrm{m}
			\]
			
			\item
			\[
			20\,\mathrm{nm}
			=
			\underline{\hspace{3cm}}\ \mathrm{m}
			\]
			
			\item
			\[
			5,0\times10^{-8}\,\mathrm{m}
			=
			\underline{\hspace{3cm}}\ \mathrm{nm}
			\]
			
			\item Uma nanopartícula possui um diâmetro de \(40\,\mathrm{nm}\).
			Quantas dessas partículas, colocadas lado a lado, seriam necessárias
			para perfazer aproximadamente \(1\,\mathrm{\mu m}\)?
			
			\linhas{2}
			
		\end{enumerate}
		
	\end{exercicio}
	
	% =========================================================
	
	\begin{exercicio}{Ordenar o mundo invisível}
		
		Considera os seguintes valores aproximados:
		
		\begin{center}
			\begin{tabular}{lc}
				\toprule
				\textbf{Objeto} & \textbf{Dimensão aproximada}\\
				\midrule
				Átomo & \(0,1\,\mathrm{nm}\)\\
				Nanopartícula de ouro & \(20\,\mathrm{nm}\)\\
				Vírus & \(100\,\mathrm{nm}\)\\
				Bactéria & \(2000\,\mathrm{nm}\)\\
				Espessura de um cabelo & \(80000\,\mathrm{nm}\)\\
				\bottomrule
			\end{tabular}
		\end{center}
		
		\begin{enumerate}[label=\alph*)]
			\item Ordena os objetos por ordem crescente de dimensão.
			\linhas{2}
			
			\item Quantas vezes a espessura aproximada de um cabelo é superior
			ao diâmetro da nanopartícula de ouro?
			\linhas{2}
			
			\item Um vírus pode ser considerado uma nanopartícula apenas com
			base na sua dimensão? Discute brevemente.
			\linhas{2}
		\end{enumerate}
		
	\end{exercicio}
	
	% =========================================================
	% ÁREA/VOLUME
	% =========================================================
	
	\begin{stem}
		
		Uma das razões pelas quais as nanopartículas são tão interessantes é a sua
		elevada \textbf{relação entre área superficial e volume}.
		
		Para uma partícula esférica de diâmetro \(d\):
		
		\[
		\frac{A}{V}=\frac{6}{d}
		\]
		
		À medida que o tamanho da partícula diminui, a relação \(A/V\) aumenta.
		
		Isto significa que uma maior fração dos átomos se encontra próxima da
		superfície, podendo tornar o material mais reativo.
		
	\end{stem}
	
	% =========================================================
	% GRÁFICO A/V
	% =========================================================
	
	\begin{figure}[H]
		\centering
		
		\begin{tikzpicture}
			
			\begin{axis}[
				width=11cm,
				height=6.5cm,
				xlabel={Diâmetro da partícula / nm},
				ylabel={Relação \(A/V\) / \(\mathrm{nm^{-1}}\)},
				xmin=0,
				xmax=110,
				ymin=0,
				ymax=0.7,
				grid=major,
				grid style={dashed,eugray},
				thick,
				mark=*,
				axis line style={eudark}
				]
				
				\addplot[
				very thick,
				mark=*,
				mark size=2.5pt
				]
				coordinates{
					(10,0.60)
					(20,0.30)
					(30,0.20)
					(40,0.15)
					(50,0.12)
					(75,0.08)
					(100,0.06)
				};
				
			\end{axis}
			
		\end{tikzpicture}
		
		\caption{Variação da relação área superficial/volume com o diâmetro
			de partículas esféricas.}
	\end{figure}
	
	% =========================================================
	
	\begin{exercicio}{Interpretar o gráfico}
		
		Usando o gráfico anterior:
		
		\begin{enumerate}[label=\alph*)]
			
			\item Indica aproximadamente \(A/V\) para uma partícula de
			\(20\,\mathrm{nm}\).
			\linhas{1}
			
			\item O que acontece à relação \(A/V\) quando o diâmetro aumenta?
			\linhas{2}
			
			\item Compara partículas de \(10\,\mathrm{nm}\) e
			\(100\,\mathrm{nm}\). Qual apresenta maior relação \(A/V\)?
			\linhas{2}
			
			\item Explica por que razão partículas muito pequenas podem ser
			particularmente interessantes em processos químicos.
			\linhas{2}
			
		\end{enumerate}
		
	\end{exercicio}
	
	% =========================================================
	% DESENHO DIVISÃO CUBO
	% =========================================================
	
	\begin{exercicio}{Por que razão partículas menores têm mais superfície?}
		
		Imagina um cubo de material com aresta \(100\,\mathrm{nm}\).
		
		Este cubo é dividido em oito cubos iguais, cada um com
		aresta \(50\,\mathrm{nm}\).
		
		\begin{center}
			
			\begin{tikzpicture}[scale=0.95]
				
				% cubo grande
				\draw[very thick,eublue,fill=eulightblue]
				(0,0) rectangle (2.5,2.5);
				
				\node[below,font=\bfseries] at (1.25,-0.2)
				{1 partícula};
				
				\node at (1.25,1.25)
				{\Large \(100\,\mathrm{nm}\)};
				
				% seta
				\draw[-{Latex[length=4mm]},very thick,eudark]
				(3.0,1.25)--(4.3,1.25);
				
				\node[above,font=\small\bfseries] at (3.65,1.35)
				{dividir};
				
				% 4 representadas frontalmente
				\foreach \x in {0,1}{
					\foreach \y in {0,1}{
						\draw[
						very thick,
						eugreen,
						fill=eulightgreen
						]
						({4.8+1.25*\x},{1.25*\y})
						rectangle
						({6.05+1.25*\x},{1.25+1.25*\y});
					}
				}
				
				\node[below,font=\bfseries] at (6.05,-0.2)
				{8 partículas};
				
				\node at (6.05,1.25)
				{\Large \(50\,\mathrm{nm}\)};
				
				\node[
				align=center,
				text=eudark,
				font=\small
				] at (6.05,3.25)
				{Representação simplificada\\da fragmentação};
				
			\end{tikzpicture}
			
		\end{center}
		
		\begin{enumerate}[label=\alph*)]
			
			\item Calcula a área superficial do cubo inicial.
			
			Recorda:
			
			\[
			A_{\text{cubo}}=6a^2
			\]
			
			\linhas{2}
			
			\item Calcula a área superficial de um dos cubos de
			\(50\,\mathrm{nm}\).
			
			\linhas{2}
			
			\item Calcula a área superficial total dos oito cubos.
			
			\linhas{2}
			
			\item Compara a área superficial total antes e depois da divisão.
			
			\linhas{2}
			
			\item O volume total de material aumentou, diminuiu ou permaneceu
			constante?
			
			\linhas{2}
			
		\end{enumerate}
		
	\end{exercicio}
	
	% =========================================================
	
	\begin{exercicio}{Nanotecnologia no quotidiano}
		
		A nanotecnologia pode ser investigada em áreas como:
		
		\begin{itemize}
			\item materiais;
			\item medicina;
			\item eletrónica;
			\item produção e armazenamento de energia;
			\item tratamento de água;
			\item sensores;
			\item revestimentos de superfícies.
		\end{itemize}
		
		Escolhe \textbf{duas} destas áreas e descreve uma possível utilização da
		nanotecnologia em cada uma.
		
		\vspace{2mm}
		
		\textbf{Área 1:}
		\linhas{3}
		
		\textbf{Área 2:}
		\linhas{3}
		
	\end{exercicio}
	
	% =========================================================
	
	\begin{exercicio}{Ciência, tecnologia e pensamento crítico}
		
		Um anúncio afirma:
		
		\begin{center}
			\begin{tcolorbox}[
				width=0.82\textwidth,
				colback=eugray,
				colframe=eupurple,
				arc=2mm
				]
				\centering
				\large\itshape
				``Este produto contém nanopartículas.\\
				Por isso, é necessariamente melhor do que qualquer produto convencional.''
			\end{tcolorbox}
		\end{center}
		
		Avalia criticamente esta afirmação.
		
		Na tua resposta considera que a palavra ``nano'' descreve uma escala
		de dimensão, mas não garante, por si só, a eficácia, utilidade ou
		segurança de um produto.
		
		\linhas{5}
		
	\end{exercicio}
	
	% =========================================================
	% QUESTÃO INTEGRADORA
	% =========================================================
	
	\section*{\textcolor{euorange}{5. Desafio final}}
	
	\begin{desafio}
		
		Um investigador estuda nanopartículas constituídas pelo elemento cobre.
		
		O cobre apresenta número atómico:
		
		\[
		Z=29
		\]
		
		Considere um átomo do isótopo:
		
		\[
		{}^{63}_{29}\mathrm{Cu}
		\]
		
		\begin{enumerate}[label=\alph*)]
			
			\item Indica o número de protões.
			\linhas{1}
			
			\item Determina o número de neutrões.
			\linhas{1}
			
			\item Indica o número de eletrões num átomo neutro.
			\linhas{1}
			
			\item Se esse átomo perder dois eletrões, quantos eletrões passa a possuir?
			\linhas{1}
			
			\item Qual será a carga do ião formado?
			\linhas{1}
			
			\item Uma nanopartícula de cobre possui um diâmetro aproximado de
			\(25\,\mathrm{nm}\). Expressa esse valor em metros.
			\linhas{2}
			
			\item Outra nanopartícula possui \(100\,\mathrm{nm}\) de diâmetro.
			Qual das duas deverá apresentar maior relação área superficial/volume?
			Justifica.
			\linhas{3}
			
			\item Explica como este exemplo estabelece uma ligação entre
			\textbf{estrutura atómica}, \textbf{materiais} e
			\textbf{nanotecnologia}.
			\linhas{3}
			
		\end{enumerate}
		
	\end{desafio}
	
	% =========================================================
	% AUTOAVALIAÇÃO
	% =========================================================
	
	\section*{\textcolor{eublue}{6. Autoavaliação}}
	
	Assinala a opção que melhor descreve a tua aprendizagem.
	
	\begin{center}
		
		\begin{tabularx}{\textwidth}{
				>{\raggedright\arraybackslash}X
				c c c
			}
			\toprule
			&
			\textbf{Ainda tenho dúvidas}
			&
			\textbf{Estou a aprender}
			&
			\textbf{Consigo sozinho}
			\\
			\midrule
			
			Distingo número atómico de número de massa.
			&
			\(\square\)&\(\square\)&\(\square\)
			\\[3mm]
			
			Determino o número de protões, neutrões e eletrões.
			&
			\(\square\)&\(\square\)&\(\square\)
			\\[3mm]
			
			Interpreto símbolos como
			\({}^{A}_{Z}\mathrm{X}\).
			&
			\(\square\)&\(\square\)&\(\square\)
			\\[3mm]
			
			Calculo uma massa atómica relativa.
			&
			\(\square\)&\(\square\)&\(\square\)
			\\[3mm]
			
			Compreendo o significado da escala nanométrica.
			&
			\(\square\)&\(\square\)&\(\square\)
			\\[3mm]
			
			Relaciono o tamanho das partículas com a razão área/volume.
			&
			\(\square\)&\(\square\)&\(\square\)
			\\
			
			\bottomrule
		\end{tabularx}
		
	\end{center}
	
	\vfill
	
	\begin{center}
		\begin{tikzpicture}
			
			\draw[eublue,very thick] (0,0) -- (13,0);
			
			\node[
			text=eublue,
			font=\Large\bfseries
			] at (6.5,0.55)
			{EUFISICA exercícios};
			
			\node[
			text=eudark,
			font=\small
			] at (6.5,-0.45)
			{Física e Química com contexto, ciência e tecnologia};
			
		\end{tikzpicture}
	\end{center}
	
\end{document}